% from ega2-1-fr.tex
\oldpage[II]{114}

de $M\otimes_A(A/\mathfrak{J})$ et sont donc images de $r$ éléments
$t_i\in M=\Gamma(Y,\mathcal{E})$. Pour tout $z_j\in Z$ il y a alors un
voisinage $V_j$ de $z_j$ tel que les restrictions des $t_i$ à $V_j$
définissent un isomorphisme
$\mathcal{O}_Y^r|V_j\to\mathcal{E}|V_j$
(\hyperref[0.5.5.4-fr]{0, 5.5.4})~; le voisinage $U$ réunion des $V_j$ répond
donc à la question.

\begin{proposition}[6.1.13]
\label{II.6.1.13-fr}
Soient $g:X'\to X$ un morphisme entier de préschémas, $Y$ un préschéma normal
localement intègre, $f$ une application rationnelle de $Y$ dans $X'$ telle
que $g\circ f$ soit une application rationnelle partout définie
(\hyperref[I.7.2.1-fr]{I, 7.2.1})~; alors $f$ est partout définie.
\end{proposition}

Si $f_1$ et $f_2$ sont deux morphismes (d'ouverts denses de $Y$ dans $X'$) de
la classe $f$, il est clair que $g\circ f_1$ et $g\circ f_2$ sont des
morphismes équivalents, ce qui justifie la notation $g\circ f$ pour leur
classe d'équivalence. On rappelle aussi que lorsqu'on suppose en outre que
$Y$ est localement noethérien, l'hypothèse que $Y$ est normal entraîne déjà
que $Y$ est localement intègre (\hyperref[I.6.1.13-fr]{I, 6.1.13}).

Pour démontrer (\hyperref[II.6.1.13-fr]{6.1.13}), notons d'abord que la
question est locale sur $Y$ et on peut par suite supposer qu'il existe dans
la classe de $g\circ f$ un morphisme $h:Y\to X$. Considérons l'image
réciproque $Y'=X'_{(h)}=X'_{(Y)}$, et notons que le morphisme
$g'=g_{(Y)}:Y'\to Y$ est entier
(\hyperref[II.6.1.5-fr]{6.1.5, (iii)}). Vu la correspondance entre
applications rationnelles de $Y$ dans $X'$ et $Y$-sections rationnelles de
$Y'$ (\hyperref[I.7.1.2-fr]{I, 7.1.2}), on voit qu'on est ramené à prouver le
cas particulier de (\hyperref[II.6.1.13-fr]{6.1.13}) où $X=Y$, autrement dit
le

\begin{corollary}[6.1.14]
\label{II.6.1.14-fr}
Soient $X$ un préschéma normal localement intègre, $g:X'\to X$ un morphisme
entier, $f$ une $X$-section rationnelle de $X'$. Alors $f$ est partout
définie.
\end{corollary}

La question étant locale sur $X$, on peut supposer $X$ intègre, et $f$
s'identifie alors à un morphisme d'un ouvert $U$ de $X$ dans $X'$
(\hyperref[I.7.2.2-fr]{I, 7.2.2}) qui est une $U$-section de $g^{-1}(U)$.
Comme $g$ est séparé, $f$ est une immersion fermée de $U$ dans $g^{-1}(U)$
(\hyperref[I.5.4.6-fr]{I, 5.4.6})~; soit $Z$ le sous-préschéma fermé de
$g^{-1}(U)$ associé à $f$ (\hyperref[I.4.2.1-fr]{I, 4.2.1}), qui est
isomorphe à $U$, donc intègre~; soit $X_1$ le sous-préschéma réduit de $X'$
ayant pour espace sous-jacent l'adhérence $\overline{Z}$ de $Z$ dans $X'$
(\hyperref[I.5.2.1-fr]{I, 5.2.1})~; $Z$ est alors un sous-préschéma induit
sur un ouvert de $X_1$ (\hyperref[I.5.2.3-fr]{I, 5.2.3}) et comme il est
irréductible, il en est de même de $X_1$, qui est par suite intègre. Le
morphisme $f$ peut alors être considéré comme une $X$-section rationnelle de
$X_1$~; comme la restriction de $g$ à $X_1$ est un morphisme entier
(\hyperref[II.6.1.5-fr]{6.1.5, (i) et (ii)}), on est finalement ramené à
prouver (\hyperref[II.6.1.14-fr]{6.1.14}) dans le cas particulier où
$X'=X_1$, autrement dit, le

\begin{corollary}[6.1.15]
\label{II.6.1.15-fr}
Soient $X$ un préschéma intègre normal, $X'$ un préschéma intègre,
$g:X'\to X$ un morphisme entier. S'il existe une $X$-section rationnelle $f$
de $X'$, $g$ est un isomorphisme.
\end{corollary}

La question étant locale sur $X$, on peut supposer $X$ affine d'anneau
intègre $A$, et alors $X'$ est affine d'anneau $A'$ entier sur $A$
(\hyperref[II.6.1.4-fr]{6.1.4}) et intègre~; en outre, le raisonnement de
(\hyperref[II.6.1.14-fr]{6.1.14}) montre qu'il y a un ouvert dense dans $X$
isomorphe à un ouvert dense de $X'$, donc $A$ et $A'$ ont même corps de
fractions. D'autre part, en vertu de
(\hyperref[I.8.2.1.1-fr]{I, 8.2.1.1}) et de l'hypothèse que les
$\mathcal{O}_x$ sont intégralement clos, l'anneau $A$ est intégralement clos,
donc $A'=A$, ce qui achève la démonstration de
(\hyperref[II.6.1.13-fr]{6.1.13}).

\subsection{Morphismes quasi-finis.}
\label{subsection:II.6.2-fr}

\begin{proposition}[6.2.1]
\label{II.6.2.1-fr}
Soient $f:X\to Y$ un morphisme localement de type fini, $x$ un point de $X$.
Les conditions suivantes sont équivalentes~:

